\subsection{The normalization lemma}

In this no. and the following, k denotes a commutative field.

THEOREM 1 (Normalization Lemma). Let A be a finitely generated \(k\) -algebra and let \(\mathfrak{a}_1 \subset \mathfrak{a}_2 \subset \ldots \subset \mathfrak{a}_p\) be an increasing finite sequence of ideals of A such that \(p \geqslant 1\) and \(\mathfrak{a}_p \neq \mathbf{A}\) . There exists a finite sequence \((x_i)_{1 \leqslant i \leqslant n}\) of elements of A which are algebraically independent over \(k\) (Chapter III, § 1, no. 1) and such that:

\begin{enumerate}
\def\labelenumi{(\alph{enumi})}
\item
  A is integral over the ring \(\mathbf{B} = k[x_1, \ldots, x_n]\) .
\item
  For all \(j\) such that \(1 \leqslant j \leqslant p\) , there exists an increasing sequence \((h(j))_{1 \leqslant j \leqslant p}\) of integers such that for all \(j\) the ideal \(\mathfrak{a}_j \cap \mathbf{B}\) of \(\mathbf{B}\) is generated \(x_1, \ldots, x_{h(j)}\) .
\end{enumerate}

Note first that it is sufficient to prove the theorem when A is a polynomial algebra \(k[Y_{1},\ldots,Y_{m}]\) . For in the general case, A is isomorphic to a quotient of such an algebra \(A'\) by an ideal \(a_{0}'\) ; let \(a_{j}'\) denote the inverse image of \(a_{j}\) in \(A'\) and let \(x_{i}'(1\leqslant i\leqslant r)\) be elements of \(A'\) satisfying the conditions of the statement for the ring \(A'\) and the increasing sequence of ideals \(a_{0}'\subset a_{1}'\subset\cdots\subset a_{p}'\) . Then the images \(x_{i}\) of the \(x_{i}'\) in A for \(i>h(0)\) satisfy the desired conditions; this is obvious for condition (b) and for condition (a) this follows from §1, no. 1, Proposition 2; finally, if the \(x_{i}(h(0)+1\leqslant i\leqslant r)\) were not algebraically independent over k, there would be a non-zero polynomial

\[
\mathrm{Q} \in k [ \mathrm{X} _ {h (0) + 1}, \dots , \mathrm{X} _ {r} ]
\]

such that \(\mathrm{Q}(x_{h(0)+1}^{\prime},\ldots,x_{r}^{\prime})\in\mathfrak{a}_{0}^{\prime}\cap\mathbf{B}^{\prime}\) , writing \(B^{\prime}=k[x_{1}^{\prime},\ldots,x_{r}^{\prime}]\) ; but by hypothesis, every element of \(a_{0}^{\prime}\cap B^{\prime}\) may be written uniquely as a polynomial in the \(x_{j}^{\prime}\ (1\leqslant j\leqslant r)\) with coefficients in k, each monomial of which contains at least one of the \(x_{j}^{\prime}\) for \(1\leqslant j\leqslant h(0)\) ; we therefore obtain a contradiction, which proves our assertion.

We shall therefore suppose in the rest of the proof that \(A = k[Y_{1}, \ldots, Y_{m}]\) and we shall argue by induction on p.

\begin{enumerate}
\def\labelenumi{(\Alph{enumi})}
\tightlist
\item
  p = 1. We distinguish two cases:
\end{enumerate}

\textbf{(A1) The ideal \(\mathfrak{a}_1\) is a principal ideal generated by an element \(x_1 \notin k\) .}\par

\(x_{1} = \mathrm{P}(\mathrm{Y}_{1},\ldots ,\mathrm{Y}_{m})\) , where \(\mathbf{P}\) is a non-constant polynomial. We shall see that, for a suitable choice of integers \(r_i > 0\) , the ring A is integral over \(\mathbf{B} = k[x_1,x_2,\dots ,x_m]\) , where

\[
x _ {i} = \mathrm{Y} _ {i} - \mathrm{Y} _ {1} ^ {r _ {i}} \quad (2 \leqslant i \leqslant m).\tag{1}
\]

For this it is sufficient to choose the \(r_i\) such that \(Y_1\) is integral over \(B\) (§ 1, no. 1, Proposition 4). Now, there is the relation

\[
\mathrm{P} \left(\mathrm{Y} _ {1}, x _ {2} + \mathrm{Y} _ {1} ^ {r _ {2}}, \dots , x _ {m} + \mathrm{Y} _ {1} ^ {r _ {m}}\right) - x _ {1} = 0.\tag{2}
\]

Let \(\mathbf{P} = \sum_{\mathbf{p}} a_{\mathbf{p}} \mathbf{Y}^{\mathbf{p}}\) , where \(\mathbf{p} = (p_1, \ldots, p_m)\) , the \(p_i\) being integers \(\geqslant 0\) , \(\mathbf{Y}^{\mathbf{p}} = \mathbf{Y}_1^{p_1} \ldots \mathbf{Y}_m^{p_m}\) , the \(a_{\mathbf{p}}\) are elements \(\neq 0\) of \(k\) and at least one of the indices \(\mathbf{p}\) is distinct from \((0, \ldots, 0)\) ; relation (2) may be written

\[
\sum_ {\mathbf {p}} a _ {\mathbf {p}} \mathrm{Y} _ {1} ^ {p _ {1}} (x _ {2} + \mathrm{Y} _ {1} ^ {r _ {2}}) ^ {p _ {2}} \dots (x _ {m} + \mathrm{Y} _ {1} ^ {r _ {m}}) ^ {p _ {m}} - x _ {1} = 0.\tag{3}
\]

Let us write \(f(\mathbf{p}) = p_{1} + r_{2}p_{2} + \cdots + r_{m}p_{m}\) and suppose that the \(r_{i}\) are chosen so that the \(f(\mathbf{p})\) are distinct (it suffices for example to take \(r_{i} = h^{i}\) , where h is an integer strictly greater than all the \(p_{j}\) (Set Theory, Chapter III, §5, no. 7, Proposition 8)). Then there will be a unique system \(\mathbf{p} = (p_{1}, \ldots, p_{m})\) such that \(f(\mathbf{p})\) is maximal and relation (3) may be written

\[
a _ {\mathbf {p}} \mathrm{Y} _ {1} ^ {f (\mathbf {p})} + \sum_ {j <   f (\mathbf {p})} \mathrm{Q} _ {j} (x _ {1}, \dots , x _ {m}) \mathrm{Y} _ {1} ^ {j} = 0\tag{4}
\]

where the \(Q_{j}\) are polynomials in \(k[Y_{1},\ldots,Y_{m}]\) ; as \(a_{p} \neq 0\) is invertible in k, (4) is certainly an equation of integral dependence with coefficients in B, whence our assertion.

The field of fractions \(k(\mathrm{Y}_1, \ldots, \mathrm{Y}_m)\) of \(\mathbf{A}\) is therefore algebraic over the field of fractions \(k(x_1, \ldots, x_m)\) of \(\mathbf{B}\) , which proves (Algebra, Chapter V, § 5, no. 3, Theorem 4) that the \(x_i\) ( \(1 \leqslant i \leqslant m\) ) are algebraically independent. Moreover, \(a_1 \cap B = Bx_1\) ; for every element \(z \in a_1 \cap B\) may be written \(z = x_1z'\) where \(z' \in A \cap k(x_1, \ldots, x_m)\) ; but \(A \cap k(x_1, \ldots, x_m) = k[x_1, \ldots, x_m] = B\) since B is integrally closed (§ 1, no. 3, Corollary 2 to Proposition 13); therefore \(z' \in B\) , which completes the proof of properties (a) and (b) in this case.

\textbf{(A2) General case \((p = 1)\) .}\par

We argue by induction on \(m\) , the case \(m = 0\) being trivial. We may obviously suppose that \(a_1 \neq 0\) (otherwise we may take \(x_i = Y_i\) for \(1 \leqslant i \leqslant m\) and \(h(1) = 0\) ). Let \(x_1\) be a non-zero element of \(a_1\) ; by (A1) there exist \(t_2, \ldots, t_m\) such that \(x_1, t_2, \ldots, t_m\) are algebraically independent over \(k\) , A is integral over \(C = k[x_1, t_2, \ldots, t_m]\) and \(x_1A \cap C = x_1C\) . By the induction hypothesis there exist elements \(x_2, \ldots, x_m\) of \(k[t_2, \ldots, t_m]\) and an integer \(h\) such that \(k[t_2, \ldots, t_m]\) is integral over \(B' = k[x_2, \ldots, x_m]\) , \(x_2, \ldots, x_m\) are algebraically independent over \(k\) and the ideal \(a_1 \cap B'\) is generated by \(x_2, \ldots, x_h\) . Then C is integral over \(B = k[x_1, x_2, \ldots, x_m]\) ( \(\S\) 1, no. 1, Corollary to Proposition 5) and hence so is A ( \(\S\) 1, no. 1, Proposition 6); the same argument as in the case (A1) shows that \(x_1, \ldots, x_m\) are algebraically independent over \(k\) ; finally, as \(x_1 \in a_1\) and \(B = B'[x_1]\) , \(a_1 \cap B = Bx_1 + (a_1 \cap B')\) and, as \(a_1 \cap B'\) is generated (in \(B'\) ) by \(x_2, \ldots, x_h\) , \(a_1 \cap B\) is generated (in B) by \(x_1, x_2, \ldots, x_h\) .

\textbf{(B) Passage from p - 1 to p.}\par

Let \(t_1, \ldots, t_m\) be elements of A satisfying the conditions of the theorem for the increasing sequence of ideals \(a_1 \subset \ldots \subset a_{p-1}\) and let us write \(r = h(p - 1)\) . By (A2) there exist elements \(x_{r+1}, \ldots, x_m\) of \(k[t_{r+1}, \ldots, t_m]\) and an integer \(s\) such that

\[
\mathbf {C} = k [ t _ {r + 1}, \dots , t _ {m} ]
\]

is integral over \(\mathbf{B}' = k[x_{r+1},\ldots,x_m]\) , \(x_{r+1},\ldots,x_m\) are algebraically independent over \(k\) and the ideal \(\mathfrak{a}_p\cap\mathbf{B}'\) is generated by \(x_{r+1},\ldots,x_s\) . Writing \(x_i=t_i\) for \(i\leqslant r\) the family \((x_i)_{1\leqslant i\leqslant m}\) obtained solves the problem with \(h(p)=s\) . For A is integral over \(\mathbf{C}[t_1,\ldots,t_r]=\mathbf{C}[x_1,\ldots,x_r]\) and hence also over \(\mathbf{B}=k[x_1,\ldots,x_m]=\mathbf{B}'[x_1,\ldots,x_r]\) since C is integral over \(\mathbf{B}'\) (§ 1, no. 1, Corollary to Proposition 5 and Proposition 6); it can be shown as in the case (A1) that the \(x_i\) are algebraically independent over \(k\) . On the other hand, for \(j\leqslant p-1\) , the ideal

\[
a _ {j} \cap k [ x _ {1}, \dots , x _ {r}, t _ {r + 1}, \dots , t _ {m} ]
\]

is by hypothesis the set of polynomials in \(x_1, \ldots, x_r, t_{r+1}, \ldots, t_m\) all of whose monomials contain one of the elements \(x_1, \ldots, x_{h(j)}\) ; as \(x_{r+1}, \ldots, x_m\) are polynomials in \(t_{r+1}, \ldots, t_m\) with coefficients in \(k\) , it is seen immediately that a polynomial in \(x_1, \ldots, x_r, x_{r+1}, \ldots, x_m\) (with coefficients in \(k\) ) can belong to \(a_j\) only if all its monomials contain one of the elements \(x_1, \ldots, x_{h(j)}\) . Finally, as \(x_1, \ldots, x_r\) belong to \(a_{p-1}\) and hence also to \(a_p, a_p \cap B'[x_1, \ldots, x_r]\) consists of the polynomials in \(x_1, \ldots, x_r\) with coefficients in \(B'\) whose constant term belongs to \(a_p \cap B'\) ; this ideal is therefore generated by \(x_1, \ldots, x_r, x_{r+1}, \ldots, x_t\) .

COROLLARY 1. Let A be an integral domain and B a finitely generated A-algebra containing A as a subring. Then there exist an element \(s \neq 0\) of A and a subalgebra B' of B isomorphic to a polynomial algebra \(\mathrm{A}[\mathrm{Y}_{1}, \ldots, \mathrm{Y}_{n}]\) such that \(\mathrm{B}[s^{-1}]\) (Chapter II, §2, no. 1) is integral over \(\mathrm{B}'[s^{-1}]\) .

We write \(S = A - \{0\}\) and let \(k = S^{-1}A\) the field of fractions of \(A\) ; clearly \(S^{-1}B\) is a finitely generated \(k\) -algebra and, as it contains \(k\) by hypothesis (Chapter II, § 2, no. 4, Theorem 1), it is not reduced to 0. By Theorem 1 (applied to \(p = 1\) and \(a_1 = 0\) ) there exists therefore a finite sequence \((x_i)_{1 \leqslant i \leqslant n}\) of elements of \(S^{-1}B\) which are algebraically independent over \(k\) and such that \(S^{-1}B\) is integral over \(k[x_1, \ldots, x_n]\) . Let \((z_j)_{1 \leqslant j \leqslant m}\) be a system of generators of the A-algebra \(B\) ; in \(S^{-1}B\) each of the \(z_j/1\) satisfies an equation of integral dependence

\[
(z _ {j} / 1) ^ {q _ {j}} + \sum_ {h <   q _ {j}} \mathrm{P} _ {h j} (x _ {1}, \dots , x _ {n}) (z _ {j} / 1) ^ {h} = 0\tag{5}
\]

where the \(P_{hj}\) are polynomials in the \(x_{i}\) with coefficients in k. There exists an element \(s \neq 0\) of A such that we may write \(x_{i} = y_{i}/s\) where \(y_{i} \in B\) for \(1 \leqslant i \leqslant n\) and all the coefficients of the \(P_{hj}\) are of the form c/s where \(c \in A\) ; finally, replacing if need be, s by a product of elements of S, we may assume that in B

\[
s z _ {j} ^ {q _ {j}} + \sum_ {h <   q _ {j}} Q _ {h j} z _ {j} ^ {h} = 0\tag{6}
\]

where the \(Q_{hj}\) are polynomials in \(y_{1},\ldots,y_{n}\) with coefficients in A; if we write \(z_{j}^{\prime}=sz_{j}\) it is seen, multiplying (6) by \(s^{q_{j}-1}\) , that \(z_{j}^{\prime}\) is integral over \(B^{\prime}=A[y_{1},\ldots,y_{n}]\) . We show that the \(y_{i}\) are algebraically independent over A; if there is a relation of the form \(\sum_{p}a_{p}y_{1}^{p_{1}}\ldots y_{n}^{p_{n}}=0\) where \(a_{p}\in A\) for all p, we deduce that \(\sum_{p}a_{p}^{\prime}x_{1}^{p_{1}}\ldots x_{n}^{p_{n}}=0\) in \(S^{-1}B\) , where \(a_{p}^{\prime}=a_{p}s^{p_{1}+\cdots+p_{n}}\) in k; by hypothesis therefore \(a_{p}^{\prime}=0\) for all p, whence \(a_{p}=0\) for all p.~Moreover, in the ring \(B[s^{-1}]\) each of the \(z_{j}^{\prime}/1\) is integral over \(B^{\prime}[s^{-1}]\) (§1, no. 1, Proposition 2) and, as \(z_{j}/1=(z_{j}^{\prime}/1)(1/s)\) in \(B[s^{-1}]\) , it is seen that the \(z_{j}/1\) are integral over \(B^{\prime}[s^{-1}]\) , which completes the proof (§1, no. 1, Proposition 4).

COROLLARY 2. Let K be a field, A a subring of K and L the field of fractions of A. If K is a finitely generated A-algebra, {[}K:L{]} is finite and there exists \(a \neq 0\) in A such that \(\mathbf{L} = \mathbf{A}[a^{-1}]\) .

It follows from Corollary 1 that there exist elements \(x_{1}, \ldots, x_{n}\) of \(\mathbf{K}\) and an element \(a \neq 0\) of \(\mathbf{A}\) such that \(x_{1}, \ldots, x_{n}\) are algebraically independent over A (and therefore over L) and that K is integral over the subring \(A[x_{1},\ldots,x_{n},a^{-1}]\) . Then it follows from §2, no. 1, Lemma 2 that \(A[x_{1},\ldots,x_{n},a^{-1}]\) is a field. But the only invertible elements of a polynomial ring \(C[Y_{1},\ldots,Y_{n}]\) over an integral domain C are the invertible elements of C; applying this remark to \(C = A[a^{-1}]\) , it is seen that necessarily n = 0 and that \(A[a^{-1}]\) is a field equal to L by definition of the latter. As K is integral over L and is a finitely generated L-algebra, the degree {[}K:L{]} is finite (§1, no. 1, Proposition 4).

COROLLARY 3. Let A be an integral domain, B a finitely generated A-algebra and b an element of B such that \(zb^n \neq 0\) for all \(z \neq 0\) in A and every integer \(n > 0\) . Let \(\rho: A \to B\) be the canonical homomorphism; there exists \(a \neq 0\) in A such that, for every homomorphism \(f\) from A to an algebraically closed field L such that \(f(a) \neq 0\) , there exists a homomorphism \(g\) from B to L such that \(g(b) \neq 0\) and \(f = g \circ \rho\) .

The hypothesis on \(b\) implies that, if \(h\) is the canonical homomorphism \(x \mapsto x / 1\) of \(\mathbf{B}\) to \(\mathbf{B}[b^{-1}]\) , the homomorphism \(h \circ \rho\) of \(\mathbf{A}\) to \(\mathbf{B}[b^{-1}]\) is injective. By Corollary 1 there therefore exist an element \(a \neq 0\) of \(\mathbf{A}\) and a subring \(\mathbf{B}'\) of \(\mathbf{B}[b^{-1}]\) such that \(\mathbf{B}[b^{-1}, a^{-1}]\) is integral over \(\mathbf{B}'[a^{-1}]\) and \(\mathbf{B}'\) is isomorphic to a polynomial algebra \(\mathbf{A}[\mathbf{Y}_1, \ldots, \mathbf{Y}_n]\) . Let \(f\) be a homomorphism from \(\mathbf{A}\) to an algebraically closed field \(\mathbf{L}\) such that \(f(a) \neq 0\) ; there exists a homomorphism from \(\mathbf{A}[\mathbf{Y}_1, \ldots, \mathbf{Y}_n]\) to \(\mathbf{L}\) extending \(f\) and hence there exists a homomorphism \(f'\) from \(\mathbf{B}'\) to \(\mathbf{L}\) extending \(f\) . As \(f'(a) \neq 0\) in \(\mathbf{L}\) , there exists a homomorphism \(f''\) from \(\mathbf{B}'[a^{-1}]\) to \(\mathbf{L}\) such that

\[
f ^ {\prime \prime} (x / a ^ {n}) = f ^ {\prime} (x). (f (a)) ^ {- n}
\]

for all \(x \in \mathbf{B}'\) and all \(n > 0\) (Chapter II, §2, no. 1, Proposition 1). Finally, as \(\mathbf{B}[b^{-1}, a^{-1}]\) is integral over \(\mathbf{B}'[a^{-1}]\) , there exists a homomorphism \(f'''\) from \(\mathbf{B}[b^{-1}, a^{-1}]\) to \(\mathbf{L}\) extending \(f''\) (§2, no. 1, Corollary 4 to Theorem 1). If \(j: x \mapsto x / 1\) is the canonical homomorphism from \(\mathbf{B}\) to \(\mathbf{B}[b^{-1}, a^{-1}]\) , \(g = f''' \circ j\) solves the problem for \(j(b)\) is invertible in \(\mathbf{B}[b^{-1}, a^{-1}]\) and hence \(f'''(j(b)) \neq 0\) in \(\mathbf{L}\) .

Note that, if B is assumed to be an integral domain and A ⊆ B in Corollary 3, the hypothesis on b is equivalent to \(b \neq 0\) .

